% ==================================================================
% id:       comb-012
% topic:    קומבינטוריקה
% subtopic: הסברים קומבינטוריים לזהויות
% solution: full
% verified: no
% ==================================================================
חשבו על שיקול דומה לזה שבשאלה הקודמת (שאלה~\qref{comb-003}), והוכיחו שלכל $n\in\N$ מתקיים
\[
\binom{2n}n=\sum_{k=0}^n\binom nk^2.
\]
(הכוונה: כמה זוגות מכנסיים יש עכשיו לגפן בארון, וכמה היא צריכה לארוז?)
\begin{solbox}
\unverified
לגפן יש עכשיו $2n$ זוגות מכנסיים, $n$ ארוכים ו-$n$ קצרים, והיא צריכה לארוז $n$ מהם. מצד אחד יש $\binom{2n}n$ דרכים.
מצד שני, נחלק למקרים לפי מספר הארוכים, $k=0,\dots,n$. אם ארזה $k$ ארוכים ו-$n-k$ קצרים, יש $\binom nk\binom n{n-k}=\binom nk^2$ דרכים (לפי הסימטריה $\binom n{n-k}=\binom nk$). לכן $\ans{\binom{2n}n=\sum_{k=0}^n\binom nk^2}$.
\end{solbox}
